% =============================================================================
% Paper: Retrieval-Augmented Trajectory Generation with Differentiable
% Geographic Constraints
%
% Target venue: ICLR 2027 (abstract ~Sep 2026)
%
% This file is the working draft. Style file (iclr2027_conference.sty or
% similar) will be added closer to submission. For now we use a neutral
% article preamble so the skeleton compiles anywhere.
% =============================================================================

\documentclass[11pt]{article}

% --- packages ----------------------------------------------------------------
\usepackage[margin=1in]{geometry}
\usepackage{amsmath,amssymb,amsthm}
\usepackage{graphicx}
\usepackage{booktabs}
\usepackage{multirow}
\usepackage{xcolor}
\usepackage{hyperref}
\usepackage{cleveref}
\usepackage{algorithm}
\usepackage{algpseudocode}
\usepackage{microtype}
\usepackage{enumitem}
\usepackage{caption}
\usepackage{subcaption}
\usepackage{url}

\hypersetup{colorlinks=true, linkcolor=blue!50!black, citecolor=blue!50!black,
            urlcolor=blue!50!black}

% --- bibliography ------------------------------------------------------------
\usepackage[numbers,sort&compress]{natbib}

% --- shortcuts ---------------------------------------------------------------
\newcommand{\TODO}[1]{\textcolor{red}{\textbf{TODO:} #1}}
\newcommand{\note}[1]{\textcolor{blue}{[#1]}}
\newcommand{\ade}{\text{ADE}}
\newcommand{\fde}{\text{FDE}}
\newcommand{\sdf}{\text{SDF}}
\DeclareMathOperator*{\argmin}{arg\,min}
\DeclareMathOperator*{\argmax}{arg\,max}

% --- title -------------------------------------------------------------------
\title{Retrieval-Augmented Trajectory Generation with\\
       Differentiable Geographic Constraints}

\author{Andrei Stirbu \\
        \texttt{andreistirbu2005@gmail.com}
        \and
        \TODO{co-authors}}

\date{\today}

% =============================================================================
\begin{document}
\maketitle

% =============================================================================
\begin{abstract}
\TODO{Final abstract. Working draft from plan \S6:}

Predicting the path a vehicle will take between two known points is a central
problem in routing, logistics, and anomaly detection. We consider this problem
for maritime vessels: given a start position, an end position, and a vessel
type, generate a 128-waypoint route consistent with historical AIS traffic and
the geography of navigable water. Existing AIS forecasting methods, including
the recent TrAISformer baseline, predict the next position from a sliding-window
history of recent reports and have no notion of long-range endpoint
conditioning or land. We propose a retrieval-augmented encoder--decoder
Transformer that (i) cross-attends to $K$ historical routes at every decoder
step, preserving per-timestep shape information lost by previous mean-pool
retrieval, (ii) is trained with a differentiable signed-distance penalty over a
coastline raster, biasing the rollout away from land, and (iii) at inference
time projects on-land predictions onto a fine-grained water-cell graph, making
the output structurally water-valid. On a US-coastal corpus of 61\,K vessel
trajectories the method reduces average displacement error by 32\% relative to a
great-circle baseline and 23\% relative to a mean-pool-retrieval ablation, with
zero land-crossings at a 10\,km clearance threshold on the held-out test set.
Cross-region transfer experiments on Danish AIS data show that the retrieval
mechanism --- but not the land loss --- transfers without retraining. We
further introduce a shared trajectory-completion benchmark that unifies
next-step prediction and endpoint-conditioned generation, enabling
apples-to-apples comparison with TrAISformer and other prior work. All code,
configs, and the cleaned corpus are released.
\end{abstract}

% =============================================================================
\section{Introduction}
\label{sec:intro}
\TODO{$\sim$1 page. Outline:}
\begin{itemize}
  \item Operational AIS use cases (ETA, anomaly detection, port management,
        environmental monitoring) want \emph{full routes} between known
        endpoints, not next-step forecasts.
  \item Prior AIS work is dominated by sliding-window next-step models
        \citep{traisformer2024,geotracknet2018}. Even when they extend to
        multi-step rollouts, they (a) compound error, (b) ignore geographic
        validity, and (c) require histories that may not be available at query
        time.
  \item We reframe the task as endpoint-conditioned route generation and
        introduce a retrieval-augmented Transformer with a differentiable
        geographic loss. We give the first head-to-head comparison with
        TrAISformer on a shared protocol.
  \item Contributions (three sharp claims):
        \begin{enumerate}[leftmargin=*,itemsep=0pt]
          \item A retrieval-augmented endpoint-conditioned Transformer for
                full-route generation with \emph{per-timestep} retrieval.
                Per-timestep retrieval preserves shape information that
                mean-pool retrieval destroys; we quantify a 23\% ADE reduction
                on long routes.
          \item A symmetric soft / hard geographic constraint: a differentiable
                signed-distance penalty during training and a discrete
                water-graph projection at inference. After
                constraint-consistent corpus construction, both terms can be
                applied without measurable ADE penalty, achieving 0\%
                land-crossings at a 10\,km clearance threshold.
          \item A unified trajectory-completion benchmark (the $K$-anchor
                protocol) that subsumes both next-step prediction and
                endpoint-conditioned generation, enabling direct comparison
                with TrAISformer and future methods.
        \end{enumerate}
  \item Headline numbers: ADE $5930 \pm 372$\,m val, $7637$\,m test on a US
        corpus of 61\,k trajectories; 0\% land-crossings; cross-region
        transfer to Danish AIS.
  \item Roadmap.
\end{itemize}

% =============================================================================
\section{Related Work}
\label{sec:related}
\TODO{$\sim$1 page. Five-paragraph structure from plan \S4; end with a
positioning paragraph that says how this work differs from each category.}

\paragraph{AIS trajectory prediction.}
\TODO{TrAISformer, GeoTrackNet, Capobianco, Forti, etc.}

\paragraph{Transformer trajectory forecasting in other domains.}
\TODO{Trajectron++, AgentFormer, HiVT, MTR.}

\paragraph{Sequence-to-sequence and graph-based route prediction.}
\TODO{Pointer Networks, NeuRoute, diffusion route planning.}

\paragraph{Constrained, physics-informed, and geography-aware NN.}
\TODO{PINN, DeepSDF, Trajectron++ constraints, safe RL.}

\paragraph{Retrieval-augmented generation.}
\TODO{RETRO, RAG, kNN-LM, retrieval for motion.}

\paragraph{Positioning.}
\TODO{One paragraph: how this paper differs from each of the five categories.}

% =============================================================================
\section{Task and Notation}
\label{sec:task}
\TODO{$\sim$\onehalf{} page.}

\subsection{Endpoint-conditioned route generation}
Given a start position $\mathbf{p}_0 \in \mathbb{R}^2$, an end position
$\mathbf{p}_T \in \mathbb{R}^2$, and a vessel-type label $v \in \mathcal{V}$,
generate a sequence of $T=128$ waypoints
$\mathbf{p}_{0:T} = (\mathbf{p}_0, \mathbf{p}_1, \dots, \mathbf{p}_T)$
\TODO{define formally; specify arc-length-uniform parameterization}.

\subsection{The $K$-anchor unified protocol}
\TODO{Define the protocol that subsumes (start, end) generation ($K=2$),
TrAISformer's history task ($K \approx 18$ at one end of the route), and
intermediate-anchor completion.}

\subsection{Metrics}
\TODO{ADE, FDE, normalized ADE, length-bucketed ADE, land\%@10\,km,
cross\%@10\,km, Frechet, inference latency.}

% =============================================================================
\section{Method}
\label{sec:method}
\TODO{$\sim$2 pages. Equations, not just prose.}

\subsection{Retrieval-augmented encoder}
\TODO{Formal definition of PerStepRouteEncoder. Encoder memory =
$3 + K \cdot t_{\text{retr}}$ tokens.}

\subsection{Decoder with windowed causal attention}
\TODO{$k_{\text{past}}=32$ window; cross-attention to encoder memory; output
head over normalized $(\Delta\text{lon}, \Delta\text{lat})$.}

\subsection{Training objective}
\begin{equation}
  \mathcal{L} = \mathcal{L}_\Delta
              + \lambda_{\text{end}}\,\mathcal{L}_{\text{end}}
              + \lambda_{\text{smooth}}\,\mathcal{L}_{\text{smooth}}
              + \lambda_{\text{land}}\,\mathcal{L}_{\text{land}}.
\end{equation}
\TODO{Define each term. Emphasize that $\mathcal{L}_{\text{land}}$ uses a
differentiable bilinear sample on the SDF raster.}

\subsection{Inference and post-processing}
\TODO{Autoregressive rollout with optional per-step hard land projection;
linear endpoint correction; WaterRouter snap-only on the 0.005\degree{} water
graph.}

\subsection{Diffusion variant (TrajDiff)}
\TODO{v13 architecture: whole-sequence denoising diffusion Transformer with
classifier-free guidance and score-based land guidance during DDIM sampling.
See \cref{sec:v13-details} for full specification.}

% =============================================================================
\section{Experiments}
\label{sec:experiments}
\TODO{$\sim$2.5 pages.}

\subsection{Datasets and protocol}
\TODO{US-coastal corpus (61\,k); Danish DMA (\TODO{n}); cleaning, MMSI-grouped
split, KNN cache.}

\subsection{Main results}
\TODO{Headline table: v10+router vs all baselines on US val and test, with
bootstrap CIs in the main table.}

\subsection{Comparison with TrAISformer}
\TODO{Three protocols from plan \S3; ADE-vs-K curves under the unified
$K$-anchor protocol.}

\subsection{Ablations}
\TODO{Per-timestep vs mean-pool retrieval; $K$ sweep; $t_{\text{retr}}$ sweep;
$k_{\text{past}}$ sweep; vessel-type drop; no-land-loss; no-router;
no-land-loss + no-router (the load-bearing ablation); hard-only vs soft-only
vs both.}

\subsection{Cross-region transfer}
\TODO{Intra-bbox (Atlantic $\leftrightarrow$ Gulf); US $\leftrightarrow$
Denmark.}

\subsection{Statistical analysis}
\TODO{Paired bootstrap 95\% CI; Bonferroni / Holm correction across baselines;
per-vessel-type breakdown; mean and median ADE across seeds.}

\subsection{Qualitative analysis}
\TODO{4-route headline figure; pre/post-snap on a peninsula; retrieval-bank
visualization.}

% =============================================================================
\section{Discussion and Limitations}
\label{sec:discussion}
\TODO{$\sim$\onehalf{} page. Pre-empt reviewer objections (see plan \S2):
post-processor doing the work; geographic bias; vessel-type fairness;
diffusion alternative.}

% =============================================================================
\section{Conclusion}
\label{sec:conclusion}
\TODO{$\sim$\onequarter{} page.}

% =============================================================================
\section*{Reproducibility statement}
\TODO{Code, configs, cleaned corpus, evaluation harness all released. Anonymized
URL placeholder for submission.}

% =============================================================================
\bibliographystyle{plainnat}
\bibliography{bib/refs}

% =============================================================================
\appendix

\section{Full architectural ladder (v4 / v5 / v9 / v10)}
\label{sec:appendix-ladder}
\TODO{Move the school-report's full ablation ladder here. Source tables:
\texttt{report/tables/full\_ladder.tex}, \texttt{tier1\_eval.csv}.}

\section{Hyperparameters and training details}
\label{sec:appendix-hp}
\TODO{Full table of HPs for v10, TrAISformer-on-US, v13.}

\section{TrajDiff (v13) full specification}
\label{sec:v13-details}
\TODO{Pull from \texttt{docs/v13\_proposal.md}; refine against actual
implementation.}

\section{Additional qualitative examples}
\label{sec:appendix-qual}
\TODO{20--40 routes with all baselines overlaid, organized by length bucket.}

\section{Failure-mode analysis}
\label{sec:appendix-failures}
\TODO{Systematic categorization of when v10+router underperforms (narrow
straits, port approaches, sparse retrieval bank, etc.).}

\end{document}
